%%%%%%%%%%%%%%%%%%%%%%%%%%%%%%%%%%%%%%%%%%%%%%%%%%%%%%%%%%%%%%%%%%%%%%%%%%%%%%%%
% Sound Similarity in YM2612 FM Presets --- TISMIR scaffold
%
% Template note: this manuscript uses the official lightweight TISMIR style
% (tismir.sty from https://github.com/ismir/paper_templates_TISMIR_old).
% For camera-ready typesetting closer to the published format, replace with
% the newer template from https://github.com/ismir/paper_templates_TISMIR_new
% (or Overleaf snip). Either official template is accepted by TISMIR.
%%%%%%%%%%%%%%%%%%%%%%%%%%%%%%%%%%%%%%%%%%%%%%%%%%%%%%%%%%%%%%%%%%%%%%%%%%%%%%%%

\documentclass{article}

\usepackage[utf8]{inputenc}
% Remove the `review` option for the camera-ready version.
\usepackage[review]{tismir}
\usepackage{amsmath}
\usepackage{booktabs}
\usepackage{graphicx}
\usepackage{hyperref}
\usepackage{url}
\usepackage{xcolor}

% Figures and tables produced by: research/ make figures  (dafm figures)
\graphicspath{{figures/}}

\title{Sound Similarity in a Large Corpus of YM2612 FM Presets:
       From Register Space to Perceptual Distance}
%
\author{%
Author Name\thanks{Affiliation, Address}%
}

\date{}

\authorref{Author,~N.}
\authorshort{Author}
% \titleshort{Sound Similarity in YM2612 FM Presets}
% \type{research}

\begin{document}

\twocolumn[{%
\maketitleblock
\begin{abstract}
% Max ~250 words. Structure only --- no invented results.
Frequency modulation (FM) synthesis on Yamaha FM chips underlies a large body
of video-game and electronic music, yet systematic evidence on how register-level
parameter changes relate to perceived timbre remains sparse.
This paper studies sound similarity over a large corpus of real-world YM2612
presets extracted from Mega Drive / Genesis games.
\textbf{Scope boundary:} the analysis starts from scratch; the only inheritance
from prior repository work is the source instrument CSV.
We render presets under a fixed protocol, extract audio descriptors and embedding
spaces, and evaluate seven experiments (E1--E7) addressing five research
questions (RQ1--RQ5).
E1 preregisters a falsifiable hypothesis from Chowning's FM theory: at equal
parametric distance, modulator total-level (TL) changes should induce larger
perceptual distance than carrier TL changes.
Hardware validation notes the instrument's 8\,MHz master clock versus the
console's $\approx$7.67\,MHz clock.
% TODO: insert quantitative summary from artifacts/results when available.
\end{abstract}
}] % end twocolumn preamble

%%%%%%%%%%%%%%%%%%%%%%%%%%%%%%%%%%%%%%%%%%%%%%%%%%%%%%%%%%%%%%%%%%%%%%%%%%%%%%%%
\section{Introduction}
\label{sec:intro}

Yamaha four-operator FM (as in the YM2612 / OPN2) maps a compact register set
to a wide range of spectra via operator routing (algorithms), ratios, envelopes
and total levels~\citep{chowning1973}.
Game-music presets provide a culturally grounded, large-scale sample of how
practitioners actually programmed the chip.
Despite progress on synthesizer programming and sound matching
\cite{yeeking2018,chen2022sound2synth,levaillant2024spinvae2},
there is still limited work that (i)~starts from a large \emph{in-the-wild}
chip-FM corpus, (ii)~separates parameter-space geometry from audio-space
geometry, and (iii)~tests theory-driven predictions such as modulator versus
carrier sensitivity.

\paragraph{Scope.}
This study does \emph{not} reuse prior cluster labels, feature choices or
hyperparameters from earlier exploratory work in the same repository.
The sole inherited artifact is
\texttt{data/processed/all\_instruments\_final.csv}.
All representations, distances and evaluation protocols are defined and
justified here.

\paragraph{Research questions.}
\begin{description}
  \item[RQ1] At equal parametric distance, do modulator TL changes yield larger
    audio/perceptual distance than carrier TL changes, as predicted by
    Chowning~\cite{chowning1973}?
  \item[RQ2] Which effective YM2612 parameters most strongly control timbral
    change under controlled perturbations?
  \item[RQ3] Do audio descriptor / embedding spaces recover known perceptual
    dimensions of timbre for FM (cf.\ McAdams et~al.~\cite{mcadams1995})?
  \item[RQ4] How do neighbourhoods in parameter space and in audio space
    relate (rank correlation; MFCC-based distances)?
  \item[RQ5] Do automatic distances agree with human odd-one-out judgments?
\end{description}

\paragraph{Contributions (claimed).}
Corpus curation and rendering protocol; multi-condition similarity spaces;
preregistered E1 and evaluation suite E1--E7 with clustered inference.
Morphing / PCA macro-controls are discussed only as future outlook
(Section~\ref{sec:outlook}) and are \emph{not} claimed as contributions.

%%%%%%%%%%%%%%%%%%%%%%%%%%%%%%%%%%%%%%%%%%%%%%%%%%%%%%%%%%%%%%%%%%%%%%%%%%%%%%%%
\section{Related Work}
\label{sec:related}

\paragraph{FM synthesis and programming.}
Chowning's formulation~\cite{chowning1973} remains the theoretical backbone.
Classical optimization for FM tone matching includes genetic and related
search~\citep{horner1993,dahlstedt2001,dahlstedt2007}.
More recent interactive and evolutionary systems continue this
line~\citep{krekovic2019}.

\paragraph{Automatic synthesizer programming and sound matching.}
Deep and hybrid approaches estimate or search parameters for VSTs and FM
instruments~\citep{yeeking2018,chen2022sound2synth,yang2023whitebox,masuda2023,roma2024}.
Latent preset models for DX7-class FM include SPINVAE /
SPINVAE-2~\citep{levaillant2020nime,levaillant2021dafx,levaillant2024spinvae2}.
Differentiable FM and DDSP-style pipelines offer alternative
routes~\citep{caspe2022ddx7,hayes2024ddsp}.
Neural proxies for black-box synths provide perceptually informed preset
embeddings~\citep{combes2025jaes}.

\paragraph{Timbre descriptors and corpora.}
The Timbre Toolbox~\citep{peeters2011} and related MDS /
perceptual studies~\citep{mcadams1995,terasawa2005,schubert2006}
motivate descriptor choices.
Chip and game-music datasets include NES-MDB~\citep{donahue2018} and
YM2413-MDB~\citep{choi2022}.
Large synthetic preset resources such as synth1B1~\citep{turian2021}
complement hand-authored banks.
Visual exploration tools for audio collections include
AudioStellar~\citep{garber2021}.
Phase vocoder and related DSP tools support analysis pipelines
\cite{driedger2014,caceres2007}.

\paragraph{Differentiation from VibeFM.}
VibeFM~\citep{rau2026vibefm} is a NIME interface for \emph{sampling and visually
exploring} FM parameter combinations, ordering glyphs by audio similarity to
support serendipitous discovery during composition.
Our work differs in goal and evidence standard: we analyse a large corpus of
\emph{existing game presets} on the YM2612, preregister theory-driven tests
(E1), and report a multi-experiment evaluation (E1--E7) aimed at MIR claims
about parameter--timbre geometry rather than an interactive composition tool.
The two lines are complementary: VibeFM targets creative workflow; we target
measurable structure in a culturally grounded chip-FM corpus.

%%%%%%%%%%%%%%%%%%%%%%%%%%%%%%%%%%%%%%%%%%%%%%%%%%%%%%%%%%%%%%%%%%%%%%%%%%%%%%%%
\section{Corpus}
\label{sec:corpus}

% TODO: insert n_presets, n_games, quarantine rates from artifacts/corpus_v1
We start from the instrument CSV
(\texttt{all\_instruments\_final.csv}) and rebuild a research corpus from
scratch: validation/quarantine rules, configuration hashing, timbral-core
grouping and TL-variant groups used by E1.
Splits and seeds are fixed in configuration
(\texttt{research/configs/}).
% TODO: insert corpus summary table from tables/ when generated.

%%%%%%%%%%%%%%%%%%%%%%%%%%%%%%%%%%%%%%%%%%%%%%%%%%%%%%%%%%%%%%%%%%%%%%%%%%%%%%%%
\section{Method}
\label{sec:method}

\subsection{Rendering}
\label{sec:render}
Presets are rendered through a YM2612 emulation core under a fixed note,
velocity, duration and gain protocol
(\texttt{research/configs/render.yaml}).
Audio is stored as FLAC with a render manifest for provenance.

\paragraph{Hardware clock note.}
Canonical Mega Drive / Genesis YM2612 clocking is approximately
7.67\,MHz (derived from the NTSC colour subcarrier chain).
The hardware instrument used for validation runs at 8\,MHz (a standard crystal).
For identical registers, output frequency scales linearly with master clock, so
an 8\,MHz instrument plays $\approx$4.3\% sharp relative to 7.67\,MHz
($\approx$+73\,cents), with envelope times shortened by the same factor.
We treat this as an analytically invertible resampling identity in the digital
domain and report residuals attributable to analogue / DAC behaviour separately
(\texttt{research/configs/hardware\_validation.yaml}).

\subsection{Features}
\label{sec:features}
We extract hand-crafted descriptors (Timbre Toolbox-inspired and loudness-aware
paths) and optional learned embeddings under versioned configs
(\texttt{research/configs/features.yaml}).
% TODO: insert feature inventory from artifacts when frozen.

\subsection{Similarity spaces}
\label{sec:spaces}
We construct multiple conditions (parameter baseline after inertness masking;
descriptor spaces; embedding spaces) with declared metrics and top-$k$ probes.
Full pairwise matrices are never materialised for $\sim$38k configurations;
rank correlations use stratified blocked sampling (E4).

%%%%%%%%%%%%%%%%%%%%%%%%%%%%%%%%%%%%%%%%%%%%%%%%%%%%%%%%%%%%%%%%%%%%%%%%%%%%%%%%
\section{Evaluation}
\label{sec:eval}

Clustered inference (within timbral core), Benjamini--Hochberg FDR and
bootstrap effect-size CIs are declared in
\texttt{research/configs/eval.yaml}.
No numerical results are reported in this scaffold.

\subsection{E1 --- Modulator vs.\ carrier TL (Chowning)}
\label{sec:e1}
Preregistered hypothesis: at equal parametric distance, pairs differing only in
modulator TL show significantly larger perceptual/audio distance than pairs
differing only in carrier TL, because modulator TL controls modulation index
(spectral bandwidth) whereas carrier TL is essentially amplitude
\cite{chowning1973}.
Primary loudness path and null-result protocol are fixed before analysis.

\subsection{E2 --- Perturbation sensitivity}
\label{sec:e2}
Budgeted one-parameter perturbations of effective registers; sensitivity curves
per parameter family.

\subsection{E3 --- McAdams-style dimensions}
\label{sec:e3}
Canonical correlation / alignment of principal audio-space axes with attack
time, spectral centroid and spectral flux~\citep{mcadams1995,peeters2011}.

\subsection{E4 --- Rank correlation across spaces}
\label{sec:e4}
Spearman correlation between parameter-space and audio-space distances on
stratified pair samples (no full matrix).

\subsection{E5 --- MFCCD anchors}
\label{sec:e5}
MFCC distance distributions with published Sound2Synth-style thresholds as
external anchors~\citep{chen2022sound2synth}.

\subsection{E6 --- Listening test}
\label{sec:e6}
Triadic odd-one-out; if infeasible, RQ5 is removed rather than left unanswered.

\subsection{E7 --- Game footprint (secondary)}
\label{sec:e7}
Same-game retrieval metrics, aggregated per game to avoid size bias; declared
confounded (OSTs intentionally mix timbres).

%%%%%%%%%%%%%%%%%%%%%%%%%%%%%%%%%%%%%%%%%%%%%%%%%%%%%%%%%%%%%%%%%%%%%%%%%%%%%%%%
\section{Results}
\label{sec:results}

% Do NOT invent numbers. Fill from research/artifacts/results after evaluation.
TODO: insert from \texttt{artifacts/results} (E1--E7 tables and figures).
Placeholder figure slots:

\begin{itemize}
  \item Figure~\ref{fig:e1} --- E1 modulator vs.\ carrier TL separation
  \item Figure~\ref{fig:e2} --- E2 perturbation sensitivity
  \item Figure~\ref{fig:spaces} --- space neighbourhood / embedding overview
\end{itemize}

% Example include (uncomment when files exist):
% \begin{figure}[t]
%   \centering
%   \includegraphics[width=\columnwidth]{e1_tl_carrier_vs_modulator.pdf}
%   \caption{TODO: caption from artifacts/results.}
%   \label{fig:e1}
% \end{figure}

%%%%%%%%%%%%%%%%%%%%%%%%%%%%%%%%%%%%%%%%%%%%%%%%%%%%%%%%%%%%%%%%%%%%%%%%%%%%%%%%
\section{Discussion and Limitations}
\label{sec:discussion}

Limitations include emulator--hardware residual error after clock alignment;
descriptor vs.\ embedding disagreements; confounding in E7; and dependence of
RQ5 on the listening study.
% TODO: discuss actual outcomes without overclaiming.

%%%%%%%%%%%%%%%%%%%%%%%%%%%%%%%%%%%%%%%%%%%%%%%%%%%%%%%%%%%%%%%%%%%%%%%%%%%%%%%%
\section{Future Outlook}
\label{sec:outlook}

PCA-based morph / macro-controls over effective parameters are a natural
\emph{downstream} application of a validated similarity space (Capacity~A in
the research package).
They are \textbf{not} claimed as contributions of this paper; any morph
results belong to follow-up systems work, not to the MIR evaluation here.

%%%%%%%%%%%%%%%%%%%%%%%%%%%%%%%%%%%%%%%%%%%%%%%%%%%%%%%%%%%%%%%%%%%%%%%%%%%%%%%%
\section{Conclusion}
\label{sec:conclusion}

We outlined a from-scratch study of sound similarity for YM2612 game presets,
centred on RQ1--RQ5 and experiments E1--E7, with a Chowning-grounded E1
hypothesis and an explicit 8\,MHz vs.\ 7.67\,MHz hardware-clock note.
% TODO: one-paragraph empirical takeaway once results exist.

%%%%%%%%%%%%%%%%%%%%%%%%%%%%%%%%%%%%%%%%%%%%%%%%%%%%%%%%%%%%%%%%%%%%%%%%%%%%%%%%
% Bibliography
%%%%%%%%%%%%%%%%%%%%%%%%%%%%%%%%%%%%%%%%%%%%%%%%%%%%%%%%%%%%%%%%%%%%%%%%%%%%%%%%
% Style comes from tismir.sty (natbib + apalike). Do not override with plain.
\bibliography{references}

\end{document}
